\section{Special Case: Relation to Least Slack Time First}
\label{app:lstf-analysis}

We now connect the auction mechanism to least slack time first (LSTF) in the deterministic stationary-goal setting from Appendix~\ref{app:deterministic-path-planning}.
Recall that each goal \(g_i\) is stationary, the agent moves deterministically, and \(d(x,g_i)\) is the minimum travel time from location \(x\) to goal \(g_i\).
At location \(x\) and time \(t\), define the slack of goal \(i\) by
\[
  \slack_i(x,t)\coloneqq d_i-t-d(x,g_i).
\]
Once \(g_i\) is reached, it disappears from the map and we write \(\slack_i(x,t)=\bot\).
The least slack time first rule schedules a goal with minimum nonnegative slack~\citep{brown2020optimal}.
Formally, at location \(x\) and time \(t\), LSTF selects any goal
\[
  i_{\mathrm{LSTF}}\in
  \arg\min_{i:\slack_i(x,t)\geq 0}
  \slack_i(x,t),
\]
with ties broken arbitrarily.
Thus LSTF gives control to the goal whose deadline is closest after accounting for travel time.

We show that, when one goal is critically urgent and all other goals are safely non-urgent, the auction selects the same goal as LSTF.
At each bidding round, every goal \(i\) is represented by a local policy that proposes an action and submits a nonnegative bid.
The highest bidder controls the agent for the next \(\tau\in\mathbb{N}_{>0}\) time steps, with ties broken uniformly at random.
Bids are bounded by \(\beta\), and \(\rho>0\) is the bid penalty coefficient.
For \poorman, bids are integer values in \([0;\beta]\); for \allpay, bids are real values in \([0,\beta]\).
We analyze two payment rules: in \poorman, only the winning policy pays \(\rho b_i\), while in \allpay, every policy pays \(\rho b_i\).

In the remainder of this appendix, \(s\in L\) denotes the current agent location and \(t\) denotes the current time.
Suppose \(\pi_i\) is the local policy for the \(i\)-th goal, mapping each pair \((s,t)\) to a bid \(b_i\) and an action \(a_i\).
We write \(\pi_i^b\colon (s,t)\mapsto b_i\) and \(\pi_i^a\colon (s,t)\mapsto a_i\) for the bidding and action components.
In this deterministic setting, \(\pi_i^a\) proceeds toward its goal \(g_i\) along a shortest path, and \(T^\tau(s,\pi_i^a)\) denotes the state reached by applying \(\pi_i^a\) for \(\tau\) steps from \(s\).

For a bidding round at time \(t\), define \(V_i^{\win}(s,t,b_i)\) and \(V_i^{\lose}(s,t,b_i)\) as the value to policy \(i\) when it bids \(b_i\) and respectively wins or loses the next bidding interval.
Let \(V_i^{\cont}(s,t)\) denote the Nash equilibrium continuation value of the subgame starting from \((s,t)\).
The value functions are defined by backward induction as
\begin{align*}
  V_i^{\win}(s,t,b_i)
  &\coloneqq
  \gamma^\tau V_i^{\cont}(T^\tau(s,\pi_i^a),t+\tau)-\rho b_i,\\
  V_i^{\lose}(s,t,b_i)
  &\coloneqq
  \gamma^\tau V_i^{\cont}(T^\tau(s,\pi_j^a),t+\tau)-\rho b_i\cdot \mathbf{1}[\text{\allpay}],
\end{align*}
where \(j\) is the policy that wins control when \(i\) loses.
The net gain from winning rather than losing is
\begin{align*}
  \Delta_i^{\mathrm{net}}(s,t,b_i)
  &\coloneqq
  V_i^{\win}(s,t,b_i)-V_i^{\lose}(s,t,b_i)\\
  &=
  \Delta_i^{\mathrm{gross}}(s,t)-\rho b_i\cdot \mathbf{1}[\text{\poorman}],
\end{align*}
where the gross gain
\[
  \Delta_i^{\mathrm{gross}}(s,t)
  \coloneqq
  \gamma^\tau\left[
    V_i^{\cont}(T^\tau(s,\pi_i^a),t+\tau)
    -
    V_i^{\cont}(T^\tau(s,\pi_j^a),t+\tau)
  \right]
\]
measures the value advantage to objective \(i\) of winning the next bidding interval.

\begin{lemma}\label{lem:app-lstf-critical-regime-gap}
  In both \poorman and \allpay settings, if policy \(i^*\) is in the critical regime, i.e., \(\slack_{i^*}(s,t)\in[0,\tau]\), and every other policy \(j\neq i^*\) is in the safe regime, i.e., \(\slack_j(s,t)>\tau\), then
  \begin{align}\label{eq:app-lstf-gross-gain-urgency}
    \Delta_{i^*}^{\mathrm{gross}}(s,t)
    -
    \max_{j\neq i^*}\Delta_j^{\mathrm{gross}}(s,t)
    > M-r_{\max}>0.
  \end{align}
\end{lemma}

\begin{proof}
  Since \(\slack_{i^*}(s,t)\leq \tau\), losing the bidding interval causes \(i^*\) to miss its deadline unless the winning policy reaches \(g_{i^*}\) incidentally during that interval.
  In the worst case represented by the losing continuation value, this incidental help does not occur.
  Hence \(V_{i^*}^{\lose}(s,t,b_{i^*})=-M-\rho b_{i^*}\cdot\mathbf{1}[\text{\allpay}]\).
  On the other hand, if \(i^*\) wins the bidding, then \(V_{i^*}^{\win}(s,t,b_{i^*})\geq -\rho b_{i^*}\), with strict inequality if the goal is reached during the next \(\tau\) steps.
  Therefore
  \[
    \Delta_{i^*}^{\mathrm{gross}}(s,t)\geq M.
  \]

  For any \(j\neq i^*\), the safe-regime condition \(\slack_j(s,t)>\tau\) implies that losing the next bidding interval does not by itself force \(j\) to miss its deadline.
  Thus \(V_j^{\win}(s,t,b_j)\leq \gamma^{\tau+1}r_j-\rho b_j\) and \(V_j^{\lose}(s,t,b_j)\geq-\rho b_j\cdot\mathbf{1}[\text{\allpay}]\), so
  \[
    \Delta_j^{\mathrm{gross}}(s,t)
    \leq
    \gamma^{\tau+1}r_j
    < r_j
    \leq r_{\max}.
  \]
  Combining the two bounds gives
  \[
    \Delta_{i^*}^{\mathrm{gross}}(s,t)-\max_{j\neq i^*}\Delta_j^{\mathrm{gross}}(s,t)
    > M-r_{\max}.
  \]
  Since \(M>r_{\max}\), the right-hand side is positive.
\end{proof}

Lemma~\ref{lem:app-lstf-critical-regime-gap} is the only special-case ingredient needed to relate the auction mechanism to LSTF.
It says that a critically urgent goal has a strictly larger gross gain than every safe goal.
Therefore, in the \poorman setting, if \(0<\rho<M-r_{\max}\), the margin condition in Theorem~\ref{thm:poorman-favors-high-control-gain} is satisfied and the critically urgent LSTF goal wins in any pure Nash equilibrium.
In the \allpay setting, the same gross-gain ordering is the input to Theorem~\ref{thm:allpay-favors-high-control-gain}; hence the mixed equilibrium concentrates control on the critically urgent LSTF goal, with concentration increasing as the urgent goal's gross gain dominates the runner-up's gross gain.
